\documentclass[11pt]{article}
\pagenumbering{gobble}
\usepackage[skip=9pt plus1pt]{parskip}
\usepackage[hmargin=20mm, top=14mm, bottom=12mm]{geometry} % US letter; matches cv.tex
\usepackage{hyperref}

%---------------------------------------------------------------------------------------
%  Fill these in per application
%---------------------------------------------------------------------------------------
\newcommand{\company}{MongoDB}
\newcommand{\role}{Software Engineer}
\newcommand{\greeting}{To whom it may concern}

\begin{document}

%---------------------------------------------------------------------------------------
%  Letterhead (matches cv.tex)
%---------------------------------------------------------------------------------------
\parbox{0.5\textwidth}{
\huge{\textbf{Minsung Cho}}
}
\hfill
\parbox{0.5\textwidth}{
\begin{tabbing}
\hspace{3cm} \= \hspace{4cm} \= \kill
\href{mailto:minsungc1@gmail.com}{\texttt{minsungc1@gmail.com}} \\
\href{https://linkedin.com/in/minsung-c}{linkedin.com/in/minsung-c}
\end{tabbing}}

\hrule
\vspace{1.5em}

\today

\vspace{1em}

\noindent \greeting{},

I'm excited to apply for the \role{} position on the Query Integration team at \company{}.
Building query engines is what I do every day, and I'd bring a track record of shipping
query engine features that hold up under real enterprise load.

At RelationalAI, I design and implement extensions to our Datalog-based relational query
language. Working across the ``full stack'' of a query, I shape the language around customer
needs, write the optimization passes that make it fast, and build the tooling and
observability to debug all of it. Recent work includes a GQL-style graph pattern matching
extension and a procedural extension that has driven previously intractable queries down to
sub-hour runtimes. I also maintain the engine's Protocol Buffer entry point, which serves
around 100k queries a day, so I know what it takes to keep a query interface reliable at scale.

This engineering is grounded in programming languages research. I designed and published a
probabilistic programming language for optimization under uncertainty that
outperforms the state of the art by up to 10x, work that earned first place at the ACM
Student Research Competition at a top programming languages conference. At \company{}, I'm looking forward to taking a
 query
language from theory to implementation at a larger scale than I've worked at before, and I'd
love to get to know the team.

\vspace{1em}

\noindent Thank you for your consideration, \\
Minsung Cho

\end{document}
